problem:  
Let \( M^n \subset \mathbb{R}^{n+1} \) be a complete, oriented hypersurface with induced metric \( g \). Assume that \( (M^n,g,f) \) satisfies the *almost gradient Ricci soliton* equation  
\[
\mathrm{Ric} + \nabla^2 f = \lambda(x) g,
\]
where \( \lambda \) is a smooth function on \( M \). Suppose further that \( M \) has polynomial volume growth and that its second fundamental form \( A \) satisfies  
\[
\int_M \big(\|A\|^2 - \tfrac{1}{n-1} H^2\big)_+ \, e^{-f} \, d\mu_g < \infty,
\]
where \( (\cdot)_+ = \max\{\,\cdot,0\,\} \).  

**Question:**  
Does such a hypersurface necessarily coincide (up to rigid motion and scaling) with a generalized round cylinder \( \mathbb{S}^k(r) \times \mathbb{R}^{n-k} \)?  
Equivalently, can one classify all complete, noncompact, almost gradient Ricci soliton hypersurfaces in Euclidean space satisfying this *integral almost-umbilicity* condition?  

Why is it a "good" problem:  
This refined formulation removes the classical pointwise pinching rigidity and replaces it with an *integral almost-umbilicity* assumption, lying at the threshold where known rigidity theorems cease to apply. It thus occupies a critical position between geometric analysis and the theory of submanifolds: analytical control via the weighted integral condition must interplay with extrinsic curvature identities and soliton potentials. The problem probes whether infinitesimal deviations from the umbilic condition, controlled on average by the soliton weight \( e^{-f} \), already enforce global rigidity. Solving it would require developing new techniques intertwining Ricci soliton potential analysis, curvature estimates, and weighted integration identities—extending both the rigidity theory for hypersurfaces and the analytic framework of almost solitons. Hence it exemplifies a "good" problem: formally simple, conceptually profound, bridging geometric flow theory with extrinsic submanifold geometry, and genuinely open at current knowledge frontiers.